\section{Reduction to the graph estimate}\label{sec:analytic-transfer}

We now explain how \Cref{prop:graph-estimate} rules out a nonzero
multiplicative correlation. The first step produces a large sum with
ordinary divisibility indicators. We then state the estimate that permits
centering those indicators, and display the resulting contradiction.
The proof of the centering estimate occupies
\Cref{sec:analytic-centering}; the graph estimate itself will be proved
in \Cref{sec:localization}.

\subsection{Absolute-value defects and a biased sequence}

A multiplicative function satisfies $f(1)=f(1)^2$. If $f(1)=0$, then
$f(n)=f(n)f(1)=0$ for every $n$. We may therefore suppose that both
functions in \Cref{thm:elliott} take the value $1$ at $1$.
There is another case in which the conclusion follows without a graph.

\begin{lemma}\label{lem:analytic-absolute-defect}
Let $f\colon\N\to\C$ be multiplicative and $|f|\le1$. If
\[
 \sum_p\frac{1-|f(p)|}{p}=\infty,
\]
then $N^{-1}\sum_{n\le N}|f(n)|\to0$.
\end{lemma}
\begin{proof}
Write $v_p(n)$ for the exponent of $p$ in $n$. For a finite set
$\mathcal P$ of primes, multiplicativity gives
\[
 |f(n)|\le
 \prod_{\substack{p\in\mathcal P\\v_p(n)=1}}|f(p)|.
\]
The majorant is periodic modulo $\prod_{p\in\mathcal P}p^2$.
Its ordinary mean is
\[
 \prod_{p\in\mathcal P}
 \left(1-(1-|f(p)|)\left(\frac1p-\frac1{p^2}\right)\right).
\]
These products tend to zero as $\mathcal P$ increases through the primes:
the sum of the subtracted quantities diverges, whereas
$\sum_p p^{-2}<\infty$. First take the long average with $\mathcal P$
fixed, and then increase $\mathcal P$.
\end{proof}

A fixed translation affects only finitely many terms of a bounded
average. Thus, after translating by the smaller shift and ordering the
functions accordingly, it suffices to prove cancellation of
$f_1(m)f_2(m+h)$ for a fixed $h\ge1$.
By \Cref{lem:analytic-absolute-defect}, we may assume
\begin{equation}\label{eq:analytic-finite-defects}
 \sum_p\frac{1-|f_i(p)|}{p}<\infty\qquad(i=1,2).
\end{equation}
If the desired cancellation fails, there are $0<\gamma\le1$ and positive
integers $N_j\to\infty$ such that
\begin{equation}\label{eq:analytic-bias}
 \left|S(N_j)\right|\ge\gamma N_j,
 \qquad S(N)=\sum_{m\le N}f_1(m)f_2(m+h).
\end{equation}
We retain this same sequence throughout the argument.
The numbers $S(N_j)$ may have varying complex arguments.

Fix $1<\tau<2$ sufficiently close to $1$ in terms of $\gamma$, and then
fix $T$ sufficiently large in terms of $\gamma$. Choose
$\phi\in C_c^\infty(\R)$ with $0\le\phi\le1$, supported in $[-T,T]$
and equal to $1$ on $[-T/2,T/2]$. For every admissible divisor $d$, use
\begin{equation}\label{eq:analytic-cutoffs}
 w_d(x)=\phi\left(\frac{k_d(x)-v_C}{\sqrt{v_C}}\right).
\end{equation}
These cutoffs satisfy the support and residue-dependence requirements
of \Cref{prop:graph-estimate}.

\subsection{Capturing the bias in one divisor interval}

Recall that $\mathcal S=\mathcal C\cup\mathcal Z$ is the finite prime set
at scale $B$, and that a squarefree product $d$ of these primes has
weight $A^{\omega_C(d)}$. The normalizing factor $L_0$ has the exact
expansion
\begin{equation}\label{eq:raw-divisor-law}
 L_0=\sum_{\substack{d\ \mathrm{squarefree}\\p\mid d\Rightarrow p\in\mathcal S}}
       \frac{A^{\omega_C(d)}}d.
\end{equation}
The term $d=1$ is included here.

\begin{lemma}[A divisor interval carrying positive mass]
\label{lem:raw-bin}
Let $f_1,f_2\colon\N\to\C$ be $1$-bounded multiplicative functions
satisfying \eqref{eq:analytic-finite-defects} and $f_1(1)=f_2(1)=1$.
Fix $1<\tau<2$. There are constants $C_0\ge1$ and
$c>0$ such that, for every sufficiently large $B$, one can choose
$1\le H\le\exp(C_0B)$ and a family $\mathcal D$ of squarefree
$\mathcal S$-products satisfying
\begin{gather}
 H<d\le\tau H,\qquad d>1,\qquad
 \omega(d)\le J=\lceil C_0\log B\rceil,
 \qquad |f_1(d)f_2(d)|\ge\tfrac12,
 \label{eq:raw-selected-divisors}\\
 \sum_{d\in\mathcal D}\frac{A^{\omega_C(d)}}d
 \ge c\frac{L_0}{B}.
 \label{eq:raw-bin-mass}
\end{gather}
The constants are independent of $B$ and of the long averaging variable.
\end{lemma}
\begin{proof}
Normalize the summands of \eqref{eq:raw-divisor-law} to a probability
law on divisors, and denote its expectation and probability by
$\E_{\mathrm{div}}$ and $\Pbb_{\mathrm{div}}$. Each prime $p$ in band $a$
is included independently with
probability $A_a/(p+A_a)$. Prime harmonic estimates give
\[
 \E_{\mathrm{div}}\omega(d)\ll_A\log B,
 \qquad
 \E_{\mathrm{div}}\log d\le A\sum_{p\le e^B}\frac{\log p}{p}\ll_A B.
\]
Choose $C_0$ large enough that the probabilities of
$\omega(d)>\lceil C_0\log B\rceil$ and $\log d>C_0B$ are each at most
$1/8$, by Markov's inequality.

For $r_p=|f_1(p)f_2(p)|$, the inequality
$1-\prod_j r_j\le\sum_j(1-r_j)$, valid when $0\le r_j\le1$, gives
\begin{align*}
 \E_{\mathrm{div}}\bigl(1-|f_1(d)f_2(d)|\bigr)
 &\le A\sum_{p>P_0}\frac{1-|f_1(p)f_2(p)|}{p}\\
 &\le A\sum_{p>P_0}\frac{(1-|f_1(p)|)+(1-|f_2(p)|)}p=o(1).
\end{align*}
Here squarefreeness permits the use of multiplicativity.
Another application of Markov's inequality shows that the probability
of $|f_1(d)f_2(d)|<1/2$ tends to zero. Also
$\Pbb_{\mathrm{div}}(d=1)=L_0^{-1}\to0$.
Thus the remaining divisors carry at least $L_0/2$ of the unnormalized
mass for large $B$.

There are at most $1+C_0B/\log\tau$ intervals
$(\tau^r,\tau^{r+1}]$, $r\ge0$, meeting $(1,e^{C_0B}]$.
One of them carries at least $cL_0/B$ of that mass, for a fixed $c>0$.
Take its lower endpoint as $H$ and retain the divisors already satisfying
the preceding conditions.
\end{proof}

For any such divisor family, any coefficients $|a_d|\le1$, and the
cutoffs \eqref{eq:analytic-cutoffs}, define the raw and centered sums by
\begin{align}
 R_B(X)&=\sum_{x\le X}\sum_{d\in\mathcal D}
 a_d A^{\omega_C(d)}f_1(x)f_2(x+hd)
 w_d(x)w_d(x+hd)\1_{d\mid x},
 \label{eq:raw-sum}\\
 C_B(X)&=\sum_{x\le X}\sum_{d\in\mathcal D}
 a_df_1(x)f_2(x+hd)g_d(x,x+hd).
 \label{eq:analytic-centered-sum}
\end{align}
The first sum retains ordinary divisibility by all primes of $d$;
the second is the sum bounded by \Cref{prop:graph-estimate}.

\begin{lemma}[Raw lower bound]\label{lem:raw-lower-bound}
Let $f_1,f_2\colon\N\to\C$ be $1$-bounded multiplicative functions,
with $f_1(1)=f_2(1)=1$, satisfying \eqref{eq:analytic-bias} and
\eqref{eq:analytic-finite-defects}. Choose $1<\tau<2$ and $T>0$ with
$\tau-1\le\gamma/64$ and $32/T^2\le\gamma/64$, and use a cutoff
$\phi$ as in \eqref{eq:analytic-cutoffs}, equal to $1$ on
$[-T/2,T/2]$. With $\mathcal D,H$ supplied by \Cref{lem:raw-bin}, set
\[
 a_d=\overline{f_1(d)f_2(d)},\qquad X_j=HN_j.
\]
There is $c_\gamma>0$, independent of $B,j$, such that, for every
sufficiently large fixed $B$,
\begin{equation}\label{eq:raw-lower-bound}
 \liminf_{j\to\infty}\frac{|R_B(X_j)|}{X_j}
 \ge c_\gamma\frac{L_0}{B}.
\end{equation}
\end{lemma}
\begin{proof}
Put
\[
 V_{\mathcal D}=\sum_{d\in\mathcal D}A^{\omega_C(d)},\qquad
 K_{\mathcal D}=\sum_{d\in\mathcal D}
   |f_1(d)f_2(d)|^2A^{\omega_C(d)}\ge\tfrac14 V_{\mathcal D}.
\]
For a fixed $d$, write $x=dm$. Unless
$\gcd(d,m)>1$ or $\gcd(d,m+h)>1$, ordinary multiplicativity gives
\[
 a_df_1(dm)f_2(d(m+h))
   =|f_1(d)f_2(d)|^2f_1(m)f_2(m+h).
\]
The exceptional set has ordinary density at most
$2\sum_{p\mid d}p^{-1}\le2J/P_0$. At fixed $B$ there are only finitely
many $d$, so all associated residue-counting errors vanish as
$j\to\infty$.

For $p\in\mathcal C$ not dividing $d$, divisibility of $dm$ by $p$ is
equivalent to divisibility of $m$. Hence the mean and variance of
$k_d(dm)$ under uniform residues are
\[
 \mu_d=v_C-\sum_{\substack{p\mid d\\p\in\mathcal C}}\frac1p,
 \qquad
 \operatorname{Var}(k_d(dm))
 =\sum_{\substack{p\in\mathcal C\\p\nmid d}}
       \frac1p\left(1-\frac1p\right)\le v_C.
\]
The same formulas hold with $m+h$ in place of $m$. For large $B$,
$|\mu_d-v_C|\le J/P_0\le(T/4)\sqrt{v_C}$. Chebyshev's inequality and
a union bound show that the proportion on which either cutoff is not
$1$ is at most $32/T^2$. Independence of the two endpoints is not needed.

The length $X_j/d$ lies between $N_j/\tau$ and $N_j$. Comparing each
shorter sum with the sum up to $N_j$ therefore gives
\begin{equation}\label{eq:raw-comparison}
 \begin{split}
 |R_B(X_j)-K_{\mathcal D}S(N_j)|
 \le N_jV_{\mathcal D}\left(
 \tau-1+\frac{32}{T^2}+\frac{4J}{P_0}+o_{j\to\infty}(1)\right).
 \end{split}
\end{equation}
The factor $4J/P_0$ allows a difference of size at most $2$ on each
coprimality failure. The endpoint rounding errors are included in $o(1)$.

Choose $\tau-1\le\gamma/64$ and $32/T^2\le\gamma/64$, and then take
$B$ and $j$ large enough for the other two errors to be at most
$\gamma/64$ each. The reverse triangle inequality, valid for the
complex number $S(N_j)$, yields
\[
 |R_B(X_j)|\ge\frac\gamma8N_jV_{\mathcal D}.
\]
Finally,
$V_{\mathcal D}\ge H\sum_{d\in\mathcal D}A^{\omega_C(d)}/d
\ge cHL_0/B$. This proves the result with $c_\gamma=\gamma c/8$.
\end{proof}

\subsection{The centering estimate and the contradiction}

The remaining analytic task is to show that centering changes the raw
sum by less than its lower bound. Its statement does not require
\eqref{eq:analytic-finite-defects}.

\begin{proposition}[Analytic centering estimate]
\label{prop:analytic-centering}
Let $f_1,f_2\colon\N\to\C$ be multiplicative, with $|f_i|\le1$, and
suppose at least one is uniformly nonpretentious. Fix
$h\ge1$, $1<\tau<2$, $C_0\ge1$, $T>0$, and a smooth function
$0\le\phi\le1$ supported in $[-T,T]$.
For each sufficiently large $B$, let $\mathcal D,H,J$ satisfy
\eqref{eq:divisor-family}, let $|a_d|\le1$, and use
\eqref{eq:analytic-cutoffs}. Define $R_B,C_B$ by
\eqref{eq:raw-sum} and \eqref{eq:analytic-centered-sum}. Then
\begin{equation}\label{eq:analytic-centering-bound}
 \limsup_{X\to\infty}\frac{|R_B(X)-C_B(X)|}{X}
 \ll L_0\left(B^{-1-\eps}+\frac{J}{P_0}\right).
\end{equation}
The implied constant may depend on the fixed parameters and $\phi$,
but is independent of the divisor family and coefficients.
The long-variable limit is taken with $B$ fixed.
\end{proposition}

\begin{proof}[Reduction of \Cref{thm:elliott}]
We deduce \Cref{thm:elliott} from
\Cref{prop:graph-estimate,prop:analytic-centering}.
The zero-function and divergent-defect cases were settled above.
Otherwise suppose \eqref{eq:analytic-bias} holds. Choose
$\tau,T,\phi,C_0$ and, for each sufficiently large fixed $B$, the divisor
family of \Cref{lem:raw-bin}. Combining \Cref{lem:raw-lower-bound} with
the two propositions gives
\[
 c_\gamma\frac{L_0}{B}
 \le C L_0\left(B^{-1-c_*/2}+B^{-1-\eps}+\frac{J}{P_0}\right).
\]
After dividing by $L_0/B$, this reads
\begin{equation}\label{eq:analytic-contradiction}
 c_\gamma\le C\left(B^{-c_*/2}+B^{-\eps}
             +BJ\exp(-B^{1-\eps})\right),
\end{equation}
which is impossible for sufficiently large $B$.

In this comparison the functions, shift, and bias $\gamma$ are fixed
first; then $\tau,T,\phi,C_0$ are fixed. For each fixed $B$ the limit
is taken along the original sequence $N_j$, with $X_j=HN_j$.
Only after these inequalities hold does $B$ increase.
Thus the contradiction excludes every alleged biased sequence.
\end{proof}

\section{The analytic centering estimate}\label{sec:analytic-centering}

We prove \Cref{prop:analytic-centering}. Expanding the centered factors
leaves sums over integers with large prime factors. We shall bound the
Fourier multiplier of those integers and use short-interval cancellation
for whichever multiplicative function is nonpretentious. Combining short
exponential sums with a fourth-moment bound follows the centering
strategy in \cite[Appendix~C]{Pilatte26}; we prove the rough-number
estimates needed for the present weights. Throughout this section,
$\mathrm e(t)=\exp(2\pi it)$.

\subsection{Expansion and finite-prime twists}

A term other than the raw term in the expansion of the center-band
factors has $d=uw$, where $w>1$ is the product of the center primes
whose constant terms were selected. Thus $u,w$ are coprime squarefree
products, every core factor of $d$ belongs to $u$, and the coefficient
of the divisibility indicator $\1_{u\mid x}$ is
\[
 A^{\omega_C(u)}a_{uw}\frac{(-\theta)^{\omega(w)}}w.
\]
For each fixed $u$, put $c_w^{(u)}=a_{uw}(-\theta)^{\omega(w)}$ on
these admissible factorizations and set $c_w^{(u)}=0$ otherwise.
Then $|c_w^{(u)}|\le1$, and its support is independent of $x$.
Set $M=H/u$ and write $x=uz$. Whenever this support is nonempty,
its values of $w$ satisfy
\begin{equation}\label{eq:analytic-rough-range}
 M<w\le\tau M,\qquad
 \frac{P_0}{\tau}\le M\le e^{C_0B}.
\end{equation}
The lower bound follows because a nonempty term has $w>P_0$.
Every such $w$ has no prime factor below $P_0$; we call integers with
this property \emph{$P_0$-rough}.

Since $u$ and $d$ have exactly the same core factors, the first
cutoff becomes
\[
 w_{uw}(uz)=\phi\left(
 \frac{\displaystyle\sum_{\substack{p\in\mathcal C\\p\nmid u}}
              \1_{p\mid z}-v_C}{\sqrt{v_C}}\right).
\]
The second cutoff has the same expression with $z+hw$ in place of $z$.
Apart from
$\gcd(u,z(z+hw))>1$, multiplicativity gives
\[
 f_1(uz)f_2(u(z+hw))=f_1(u)f_2(u)f_1(z)f_2(z+hw).
\]
For a fixed $u,w$, the exceptional set has ordinary density at most
$2J/P_0$. Since
\begin{equation}\label{eq:analytic-sum-u}
 \sum_{M<w\le\tau M}\frac1w\ll_\tau1,
 \qquad
 \sum_{\substack{u\ \mathrm{squarefree}\\p\mid u\Rightarrow p\in\mathcal S}}
       \frac{A^{\omega_C(u)}}u=L_0,
\end{equation}
the total error, divided by $X$, has limsup $O(L_0J/P_0)$.
The case of an identically zero factor is immediate, so in this
factorization we may assume $f_1(1)=f_2(1)=1$.

For $\widehat\phi(t)=\int_{\R}\phi(s)\mathrm e(-ts)\,ds$,
Fourier inversion expresses each cutoff as an integral of constant
phases times twists
\begin{equation}\label{eq:analytic-twists}
 b_i(n)=f_i(n)\mathrm e\left(
 \frac{t_i}{\sqrt{v_C}}
 \sum_{\substack{p\in\mathcal C\\p\nmid u}}\1_{p\mid n}\right),
 \qquad t_i\in\R.
\end{equation}
For coprime integers the count in this exponent is additive. Hence
$b_i$ is multiplicative and $|b_i|\le1$. It need not be completely
multiplicative: the added factor has the same value at $p$ and $p^2$.
Its prime values agree with those of $f_i$ outside $\mathcal C$.
The two Fourier integrals have total absolute weight
$\|\widehat\phi\|_1^2$, a fixed finite constant.

It is therefore enough to show, uniformly in the twists
\eqref{eq:analytic-twists} and in $|c_w|\le1$ supported on
$P_0$-rough integers in $(M,\tau M]$, that
\begin{equation}\label{eq:analytic-rough-target}
 \limsup_{Y\to\infty}\frac1Y
 \left|\sum_{z\le Y}b_1(z)
       \sum_{M<w\le\tau M}\frac{c_w}{w}b_2(z+hw)\right|
 \ll B^{-1-\eps}.
\end{equation}
Here $Y=X/u$, and all graph parameters are fixed before $Y$ increases.

\subsection{The short-interval input}

We first justify the uniformity in the Fourier parameters. Finite
changes to prime values have only a bounded effect on the squared
distance defining nonpretentiousness.

\begin{lemma}[Stability under finitely many prime changes]
\label{lem:finite-prime-stability}
Let $f,b\colon\N\to\C$ be multiplicative and $1$-bounded. Suppose their
prime values agree outside a finite set $\mathcal P$. For every
Dirichlet character $\chi$, every real $t$, and $X\ge2$,
\begin{equation}\label{eq:analytic-distance-stability}
 \left|D(b,\chi n^{it};X)^2-D(f,\chi n^{it};X)^2\right|
 \le2\sum_{p\in\mathcal P}\frac1p.
\end{equation}
Consequently, if $f$ is uniformly nonpretentious, the same divergence
holds uniformly over all such $b$ and over every fixed finite family
of Dirichlet characters.
\end{lemma}
\begin{proof}
Outside $\mathcal P$ the summands in the squared distances coincide.
At a prime in $\mathcal P$ their difference has absolute value at most
$|b(p)-f(p)|/p\le2/p$. Sum this bound, and then take the infimum over
$|t|\le X$ and the minimum over the finite character family.
\end{proof}

For a $1$-bounded multiplicative function $b$, define
\[
 \mathcal M(b;X,Q)=
 \inf_{\substack{|t|\le X\\q\le Q,\ \chi\ ({\rm mod}\ q)}}
     D(b,\chi n^{it};X)^2.
\]
The following is the general exponential-sum theorem of
Matom\"aki, Radziwi\l\l, and Tao, in its corrected version
\cite[Theorem~1.7]{MRT15}.

\begin{theorem}[Averaged short exponential sums]\label{thm:mrt}
Let $X\ge D\ge10$ and let $b\colon\N\to\C$ be a $1$-bounded
multiplicative function. With
$Q=\min((\log X)^{1/125},(\log D)^5)$, one has
\begin{equation}\label{eq:mrt-input}
 \begin{split}
 \sup_{\alpha\in\T}\int_0^X
 \left|\sum_{y<m\le y+D}b(m)\mathrm e(\alpha m)\right|\,dy
 \ll DX\left(
 e^{-\mathcal M(b;X,Q)/20}
 +\frac{\log\log D}{\log D}
 +\frac1{(\log X)^{1/700}}\right).
 \end{split}
\end{equation}
The implied constant is absolute.
\end{theorem}

The interval endpoint convention does not affect the integral.
For fixed $B$ and $D$, the characters of moduli at most $(\log D)^5$
form a finite family. If $f_i$ is uniformly nonpretentious,
\Cref{lem:finite-prime-stability} with $\mathcal P=\mathcal C$ gives
$\mathcal M(b_i;Y,Q)\to\infty$ uniformly in the twists
\eqref{eq:analytic-twists}. It follows that
\begin{equation}\label{eq:mrt-fixed-length}
 \limsup_{Y\to\infty}\sup_{\alpha,t_i}
 \frac1{DY}\int_0^Y
 \left|\sum_{y<m\le y+D}b_i(m)\mathrm e(\alpha m)\right|\,dy
 \ll\frac{\log\log D}{\log D}.
\end{equation}
Both the prime cutoff and the allowed height in the distance are $Y$,
as in the hypothesis of \Cref{thm:elliott}. For real $Y$, apply
\Cref{thm:mrt} with $X=\lceil Y\rceil$ and enlarge the integral to
$[0,\lceil Y\rceil]$; the normalization changes by a factor tending
to $1$. Thus the hypothesis along integer scales is sufficient.
The frequency supremum in
\eqref{eq:mrt-fixed-length} is outside the integral. No bound with that
supremum inside the integral is used below.

\subsection{A rough-number Fourier multiplier}

To use \eqref{eq:mrt-fixed-length} for
\eqref{eq:analytic-rough-target}, define
\begin{equation}\label{eq:rough-multiplier}
 Q_M(\alpha)=\sum_{M<w\le\tau M}\frac{c_w}{w}\mathrm e(\alpha hw).
\end{equation}
We need its maximum and its fourth moment. The coefficients are
arbitrary; only their rough support will be used.

\begin{lemma}[Rough-number bounds]\label{lem:rough-multiplier}
Fix $C_0\ge1$, $1<\tau<2$, and an integer $h\ge1$. Let
$P_0=\exp(B^{1-\eps})$, where $\eps=10^{-4}$, and suppose
$P_0/\tau\le M\le\exp(C_0B)$. If $|c_w|\le1$ and $c_w=0$ unless
$w\in(M,\tau M]$ is $P_0$-rough, then, for all sufficiently large $B$,
\begin{align}
 \|Q_M\|_\infty&\ll B^{-1+\eps}\log B,
 \label{eq:rough-sup}\\
 \int_{\T}|Q_M(\alpha)|^4\,d\alpha
 &\ll M^{-1}B^{-4+4\eps}(\log B)^6.
 \label{eq:rough-fourth}
\end{align}
The constants are uniform in $M$ and the coefficients.
\end{lemma}
\begin{proof}
We give the sieve bounds including their counting errors. Put
\[
 s=2\lceil100\log B\rceil,\qquad z_0=P_0^{1/(100s)}.
\]
For every prime $p<z_0$, forbid $\nu_p$ residue classes, where
$0\le\nu_p\le2$. If $r(n)$ is the number of these prime conditions
satisfied by $n$, even inclusion-exclusion gives
\[
 \1_{r(n)=0}\le\sum_{j=0}^s(-1)^j\binom{r(n)}j.
\]
Indeed, for $r>0$ the sum on the right is $\binom{r-1}s\ge0$, and for
$r=0$ it is $1$.
Write $E_j((a_p))=\sum_{|\mathcal I|=j}\prod_{p\in\mathcal I}a_p$
for the elementary symmetric sum over subsets of the primes $p<z_0$.
The Chinese remainder theorem, applied on any interval $I$ of real
length $V$, now yields
\begin{equation}\label{eq:rough-brun}
 \begin{split}
 \#\{n\in I:n\text{ avoids the forbidden classes}\}
 \le V\sum_{j=0}^s(-1)^jE_j((\nu_p/p))
       +O\left(\sum_{j=0}^sE_j((\nu_p))\right).
 \end{split}
\end{equation}
The residue-counting error is independent of the position of $I$.
It is at most
\[
 (s+1)(2z_0)^s=(s+1)2^sP_0^{1/100}\le P_0^{1/5}
\]
for large $B$.

Let $\Lambda_0=\sum_{p<z_0}\nu_p/p$. Mertens' estimate gives
$\Lambda_0\le2\log\log z_0+O(1)\le3\log B$ for large $B$.
Since $E_j((\nu_p/p))\le\Lambda_0^j/j!$, the difference between the
truncated density in \eqref{eq:rough-brun} and its full Euler product
has absolute value at most
\[
 \sum_{j>s}\frac{\Lambda_0^j}{j!}
 \le\frac{(e\Lambda_0/(s+1))^{s+1}}{1-\Lambda_0/(s+2)}
 \le B^{-10}.
\]
Consequently the count in \eqref{eq:rough-brun} is at most
\begin{equation}\label{eq:rough-sieve-bound}
 V\left(\prod_{p<z_0}(1-\nu_p/p)+B^{-10}\right)+O(P_0^{1/5}).
\end{equation}

For one rough integer, take $\nu_p=1$. The product is
$O(1/\log z_0)=O(B^{-1+\eps}\log B)$. Thus the number $R$ of
$P_0$-rough integers in $(M,\tau M]$ satisfies
\begin{equation}\label{eq:rough-one-point}
 R\ll M B^{-1+\eps}\log B.
\end{equation}
Both error terms in \eqref{eq:rough-sieve-bound} are absorbed, since
$M\ge P_0/\tau$.

The one-point bound proves \eqref{eq:rough-sup}, because $w>M$.
To estimate the fourth moment, extend $a_w=c_w/w$ by zero outside its
support. Fourier orthogonality gives
\[
 \int_{\T}|Q_M(\alpha)|^4\,d\alpha
 =\sum_k\left|\sum_w a_{w+k}\overline{a_w}\right|^2.
\]
There is no factor depending on $h$: multiplication by the nonzero
integer $h$ preserves Haar measure on $\T$.
Let $R_k$ count pairs $w,w+k\in(M,\tau M]$ that are both
$P_0$-rough. Each inner absolute value is at most $R_k/M^2$.
We next bound $R_k$ for $k\ne0$; only $|k|\le2M$ can contribute.
The permitted interval for $w$ is an intersection of two intervals of
length less than $M$. Sieve out the residues $0,-k$ modulo $p<z_0$.
If $k$ is odd there are no pairs for large $B$, because both rough
integers must be odd. If $k$ is even, the density product is
\begin{align*}
 \frac12\prod_{3\le p<z_0}(1-2/p)
 \prod_{\substack{3\le p<z_0\\p\mid k}}
       \frac{1-1/p}{1-2/p}
 &\ll\frac1{(\log z_0)^2}
       \exp\left(\sum_{p\mid k}\frac1p+O(1)\right).
\end{align*}
Here $(1-2/p)/(1-1/p)^2=1-1/(p-1)^2\le1$ and
$\log((1-1/p)/(1-2/p))=1/p+O(p^{-2})$ for $p\ge3$.
The singular factor is uniformly $O_{C_0}(\log B)$: indeed
\[
 \sum_{p\mid k}\frac1p
 \le\sum_{p\le B}\frac1p+\frac{\log|k|}{B\log B}
 \le\log\log B+O_{C_0}(1),
\]
since $\log|k|\le C_0B+\log2$. Applying
\eqref{eq:rough-sieve-bound} proves
\begin{equation}\label{eq:rough-two-point}
 R_k\ll M B^{-2+2\eps}(\log B)^3\qquad(k\ne0).
\end{equation}

There are $O(M)$ nonzero differences, so their total contribution is
$O(M^{-1}B^{-4+4\eps}(\log B)^6)$ by
\eqref{eq:rough-two-point}. The term $k=0$ is at most
$(R/M^2)^2\ll M^{-2}$, which is smaller than the same bound because
$M\ge P_0/\tau$. This proves \eqref{eq:rough-fourth}.
\end{proof}

The fourth moment localizes the frequencies at which $Q_M$ is large.
On the remaining frequencies its maximum is already small enough.
We next put the target correlation into a form where these two facts
can be combined with \eqref{eq:mrt-fixed-length}.

\subsection{Two overlapping intervals and the frequency split}

Let
\[
 D_1=\lceil M\rceil,\qquad D_2=(1+2h)D_1,
\]
and define
\[
 A_y(\alpha)=\sum_{y<n\le y+D_1}b_1(n)\mathrm e(\alpha n),
 \qquad
 B_y(-\alpha)=\sum_{y<m\le y+D_2}b_2(m)\mathrm e(-\alpha m).
\]
For each term of $Q_M$, integration in $\alpha$ imposes $m=n+hw$.
Moreover $hw<2hD_1=D_2-D_1$, so
\[
 [n-D_1,n)\subseteq[n+hw-D_2,n+hw).
\]
Thus, if $L_Y(n)=|[0,Y]\cap[n-D_1,n)|$, orthogonality gives the exact
identity
\[
 \int_0^Y\int_{\T}Q_M(\alpha)A_y(\alpha)B_y(-\alpha)
                 \,d\alpha\,dy
 =\sum_w\frac{c_w}{w}\sum_{n\ge1}b_1(n)b_2(n+hw)L_Y(n).
\]
For $D_1\le n\le Y$ one has $L_Y(n)=D_1$. The difference from
$D_1\1_{1\le n\le Y}$ is supported on $O(D_1)$ integers at the two
ends and has total absolute mass $O(D_1^2)$. Since
$\sum_{M<w\le\tau M}1/w\ll1$, we obtain
\begin{equation}\label{eq:analytic-overlap}
 \begin{split}
 &\frac1{D_1Y}\int_0^Y\int_{\T}
        Q_M(\alpha)A_y(\alpha)B_y(-\alpha)\,d\alpha\,dy\\
 &\hspace{10mm}=
 \frac1Y\sum_{n\le Y}b_1(n)
       \sum_{M<w\le\tau M}\frac{c_w}{w}b_2(n+hw)
       +O(D_1/Y).
 \end{split}
\end{equation}
The endpoint error tends to zero because $D_1$ is fixed before $Y$.

Set $t_B=B^{-1-\eps}$ and
$E_B=\{\alpha\in\T:|Q_M(\alpha)|>t_B\}$.
By \Cref{lem:rough-multiplier},
\begin{equation}\label{eq:rough-large-frequencies}
 |E_B|\ll M^{-1}B^{8\eps}(\log B)^6,
 \qquad \|Q_M\|_\infty\ll B^{-1+\eps}\log B.
\end{equation}
On $\T\setminus E_B$, Parseval and Cauchy--Schwarz give, for each $y$,
\[
 \int_{\T\setminus E_B}|Q_M A_y B_y|\,d\alpha
 \le t_B\left(\int_{\T}|A_y|^2\right)^{1/2}
             \left(\int_{\T}|B_y|^2\right)^{1/2}
 \le t_B\sqrt{D_1D_2}.
\]
After integration in $y$ and division by $D_1Y$, the contribution is
$O_h(B^{-1-\eps})$.

On $E_B$, suppose first that $f_1$ is uniformly nonpretentious.
Use $|B_y|\le D_2$, integrate first in $y$, and apply
\eqref{eq:mrt-fixed-length} to $A_y$ at each fixed frequency. The
normalized limsup is at most
\begin{equation}\label{eq:analytic-large-bound}
 O\left(\|Q_M\|_\infty\,|E_B|\,D_2
                    \frac{\log\log D_1}{\log D_1}\right).
\end{equation}
If instead $f_2$ is the nonpretentious factor, use $|A_y|\le D_1$
and apply the same estimate to $B_y$; this replaces $D_1$ by $D_2$
inside the logarithms in \eqref{eq:analytic-large-bound}.
In both cases Fubini's theorem is followed by the bound for the
already integrated exponential sum, with a supremum over its fixed
frequency. The frequency supremum has not been moved inside a
short-interval integral.

The range \eqref{eq:analytic-rough-range} implies, for $i=1,2$,
\[
 \frac{D_i}{M}\ll_h1,
 \qquad
 \frac{\log\log D_i}{\log D_i}
 \ll_{C_0,h,\tau}B^{-1+\eps}\log B.
\]
Using \eqref{eq:rough-large-frequencies} in
\eqref{eq:analytic-large-bound}, the large-frequency contribution is
\[
 O\left(B^{-2+10\eps}(\log B)^8\right)
   =o(B^{-1-\eps}),
\]
since $11\eps<1$. Together with the small-frequency bound and
\eqref{eq:analytic-overlap}, this proves
\eqref{eq:analytic-rough-target}. The argument is uniform in the
Fourier twists, since \eqref{eq:mrt-fixed-length} is uniform in them.

\begin{proof}[Completion of \Cref{prop:analytic-centering}]
Apply \eqref{eq:analytic-rough-target} with $c_w=c_w^{(u)}$ to each fixed
$u$ in the expansion at the start of this section. Dividing its contribution by
$X$ contributes the factor $A^{\omega_C(u)}/u$; the extracted factors
$f_1(u)f_2(u)$ have modulus at most $1$.
The two Fourier inversions cost at most $\|\widehat\phi\|_1^2$.
At fixed $B$ all sums over $u,w$ are finite, and the uniform estimate
therefore gives, by \eqref{eq:analytic-sum-u}, a total limsup
$O(L_0B^{-1-\eps})$. Adding the earlier coprimality error
$O(L_0J/P_0)$ proves \eqref{eq:analytic-centering-bound}.
\end{proof}

The analytic comparison required for
\eqref{eq:analytic-contradiction} is now proved. It remains to establish
the graph estimate; the following sections develop its closed-line
bound and then pass from that bound to the long-average estimate.
